% 원제: 16. 향후 계획
\chapter{향후 계획}

\section{남은 작업}

최종 보고(2026년 11월 말 예정)까지 남은 작업과 완료 기준은 표 \ref{t:plan}과 같다. 검증 결과가 보고서에 한 번에 들어가도록 검증 과제를 먼저 끝내고 최종 보고서는 마지막에 쓴다. 표의 1–5는 최종 보고서의 결론을 바꾸거나 받치는 과제로 보고서 전에 하고, ‘후속’은 결론을 바꾸지 않아 한계와 후속 과제로 넘긴다.

{\footnotesize
\begin{longtable}{L{\dimexpr 0.0702\linewidth-2\tabcolsep\relax}L{\dimexpr 0.2456\linewidth-2\tabcolsep\relax}L{\dimexpr 0.3509\linewidth-2\tabcolsep\relax}L{\dimexpr 0.3333\linewidth-2\tabcolsep\relax}}
\caption{남은 작업, 필요성과 완료 기준}\label{t:plan}\\
\toprule
\tabhead \textbf{순서} & \textbf{작업} & \textbf{필요성} & \textbf{완료 기준} \\
\midrule
\endfirsthead
\toprule
\tabhead \textbf{순서} & \textbf{작업} & \textbf{필요성} & \textbf{완료 기준} \\
\midrule
\endhead
\bottomrule
\endlastfoot
1 & 영교차 게이트 방향 반전 처리 & 지금 남은 유일한 회로 불통과다(TB20 사례 D: 반주기 첨두가 무장 문턱보다 작은 방향 반전에서 턴오프 1.01 A) & 무장하지 못한 반주기를 공칭 도통시간에 끄는 게이트로 같은 사례를 새 테스트벤치에서 재판정. 기존 TB18–TB20 사례 회귀 \\
2 & 선행기술 전수 검색 & 기여 주장이 12건 조사 범위의 결과다. 같은 조합이 나오면 기여를 다시 써야 하므로 일찍 확인한다 & IEEE Xplore·특허 데이터베이스 검색과 청구항 대조, 13.3절 갱신 \\
3 & Stage A 모듈·위상 셰딩 조합과 공차 & R2의 핵심인 Stage A를 실제 4모듈 8위상으로 경부하 운전하는 순서와, 위상별 공진주파수 공차에서의 셰딩 문턱을 확인하지 않았다 & 모듈 셰딩 + 위상 셰딩 상태기계, 공차 표본 시험, 대표 사례 회로 확인 \\
4 & SALS 분리·감발 연동 & 고장 강등 순서가 400 V 부하 분리·감발 명령에 의존하는데, 부하 승인 계층(H4)은 축약 모델에서만 확인했다 & 보호 상태기계 → SALS 명령 → 부하 응답 지연을 포함한 시험 \\
5 & 부품 정격 표 & V\textsubscript{A} 클램프, 제동 저항, 바이패스 선형 영역, 누설·필터 인덕터 포화 전류, Stage C 데드타임을 부품 사양으로 넘겨야 한다. 1–4에서 요구가 더 나올 수 있어 마지막에 한다 & 데이터시트 대조와 정격 표, 반복 고장 열 계산 \\
6 & 최종 보고서 & 연구 주장을 “Stage A 중심 다기능 허브 + 팩 등급 조건부 창 조정 모듈(R2)”과 그 성립 조건으로 정리 & 본 보고서를 뼈대로 1–5의 결과를 넣는다. H5 기각, 산업 방식 대비 조건, 회로의 불통과 사례를 함께 보고 \\
후속 & DAB 지령표 연속화 & TB13의 전이 여유 5 \%를 되찾는 최적화로, 이미 1.5배 기준을 통과한 설계의 결론을 바꾸지 않는다 & 전류 방향 연속 제약을 둔 지령표로 전이 여유 없이 회로 1.5배 통과 \\
후속 & 하드웨어 실증 & 실제 지연·잡음·열·전자기 간섭은 시뮬레이션으로 확인할 수 없다 & HIL, 축소전력 시제품 \\
\end{longtable}
}

\section{수행 순서}

검증 과제(1–5)를 11월 중순까지 끝내고, 그 결과를 넣어 최종 보고서를 11월 말까지 정리한다. 게이트 처리는 Stage A 게이트 함수의 작은 변경이고 회로 사례 하나로 판정할 수 있어 가장 먼저 한다. 선행기술 검색은 기여 주장을 바꿀 수 있어 일찍 한다. 부품 정격 표는 앞의 검증에서 요구가 더 나올 수 있어 검증 과제의 마지막에 둔다. 단계마다 대표 사례를 새 Simulink 테스트벤치로 확인하는 방식은 그대로 유지한다.
